\section{Architecture}

The pipeline is a sequence of well defined stages. A Python importer reads the
ONNX model and, using the MLIR Python bindings, builds \texttt{npu} dialect IR, a
tensor level representation whose values conceptually live in DRAM. A sequence of
passes optimizes this IR to a fixed point. The dialect conversion framework then
lowers it to the \texttt{npuisa} dialect, an instruction level representation in
which values are scratchpad resident buffers and data crosses the DRAM boundary
only through explicit DMA. A scratchpad allocator assigns each buffer a byte
offset and spills to DRAM when the working set exceeds the budget. The instruction
encoder serializes the result to a documented binary format, which a disassembler
can render and which the simulator executes.

The compiler is a small out of tree CMake project. It never builds LLVM itself; it
links against a prebuilt LLVM and MLIR at a pinned release tag. This keeps the
per change build in the seconds to minutes range after the one time toolchain
build. The driver, \texttt{npu-compile}, wraps the whole pipeline with
optimization levels and staged output so any intermediate can be inspected.
